\begin{center}
\begin{tikzpicture}[x=.9cm,y=.9cm,font=\small]
  \fill[paper] (0,0) rectangle (10.4,5.4);
  \draw[blue!55,line width=1pt] (.3,4.5) .. controls (2.8,4.2) and (5.0,4.8) .. (9.9,4.25);
  \node[text=blue!70!black,above] at (5.0,4.45) {黄河};
  \fill[dynasty!23] (1.0,2.0) -- (3.4,1.9) -- (3.7,3.7) -- (1.4,3.9) -- cycle;
  \node[align=center] at (2.35,2.9) {宗周\\关中};
  \fill[timeline!24] (5.0,2.3) -- (6.8,2.2) -- (7.0,3.8) -- (5.3,4.0) -- cycle;
  \node[align=center] at (6.0,3.1) {成周\\洛邑};
  \fill[green!18] (8.0,1.2) -- (10.0,1.3) -- (10.1,3.6) -- (8.3,3.8) -- cycle;
  \node[align=center] at (9.1,2.5) {东方封国\\与旧商区域};
  \draw[-{Stealth[length=2mm]},ultra thick,color=dynasty] (3.5,2.9) -- (4.85,3.0);
  \draw[-{Stealth[length=2mm]},ultra thick,color=timeline] (7.1,3.0) -- (7.85,2.7);
  \node[text=dynasty,above] at (4.15,3.1) {王室东向支点};
  \node[text=timeline,below] at (7.55,2.7) {安置与经营};
\end{tikzpicture}

\smallskip
\small\textit{图  宗周、成周与东方网络示意：显示双重政治重心、东方经营和封国网络的相对关系；不表示精确迁置路线或城市边界。}
\end{center}
